% e_entries.tex - Terms beginning with E
\chapter*{E}
\addcontentsline{toc}{chapter}{E}

\dictentry{EFT Switch}{/ee eff tee switch/, \textit{n.}}{
Electronic Funds Transfer Switch. A computer system that acts as a central hub connecting different banks, card networks, and payment providers together. When you use your debit card at an ATM or store, the EFT Switch routes your transaction request to your bank and brings back the approval or decline response. It handles real-time message routing, protocol translation, and ensures transactions reach the correct destination for authorization.

\example{When you withdraw cash from an ATM that belongs to a different bank than yours, an EFT Switch connects your bank to that ATM's bank to process the transaction in seconds.}

\analogy{Biological}{The EFT Switch functions like the nervous system's neural switching centers (ganglia). Just as sensory signals must be routed through neural hubs to reach the brain for processing, payment requests must pass through the EFT Switch to reach the appropriate bank. The switch performs ``address resolution''---determining which bank owns the card---similar to how the nervous system routes signals to specific brain regions. It's a high-throughput, fault-sensitive system where timing matters: delayed routing can cause transaction timeouts, just as neural delays cause impaired responses.}

\crossref{\textbf{ATM}, \textbf{Payment Gateway}, \textbf{Authorization}, \textbf{Payment Network}}
}
\indexentry{EFT Switch}

\dictentry{EMV}{/ee em vee/, \textit{abbr.}}{
Europay, Mastercard, and Visa. A global standard for smart payment cards that use an embedded chip to perform dynamic authentication for each transaction. Unlike magnetic stripe cards that contain static data (easily copied and reused forever), EMV chips generate unique cryptographic data for every transaction. The card validates the terminal, the terminal validates the card, and the system supports both offline and online risk checks. This checkpoint-based approach prevents cloning and replay attacks.

\example{When you insert your credit card into a payment terminal at a store and wait a few seconds, you're using EMV chip technology. The chip is generating a unique code for this specific transaction that cannot be reused, even if intercepted.}

\analogy{Biological}{EMV behaves like an immune checkpoint system in living cells. A magnetic stripe card is like a dead identifier---if copied once, it can be reused forever (similar to static DNA without repair mechanisms). An EMV chip, however, behaves like a living cell: it responds differently each time, checks its environment, and refuses interaction if conditions look wrong. Each EMV transaction is like a controlled biochemical reaction---time-bound, context-aware, and hard to fake outside the system. Card cloning is analogous to cancer cell impersonation, while EMV's dynamic authentication works like immune checkpoint verification. Offline risk limits mirror local tissue-level decision making, and online issuer checks function as central regulatory override. If the ``cell'' behaves abnormally, the system escalates.}

\crossref{\textbf{Chip Card}, \textbf{Magnetic Stripe}, \textbf{Tokenization}, \textbf{POS}, \textbf{Cryptogram}}

\etymology{Named after the three companies (Europay, Mastercard, Visa) that originally created this standard in the 1990s. EMV chips became widely adopted in the United States after 2015.}
}
\indexentry{EMV}

